% !TEX TS-program = lualatex
На вход интеллектуального агента должны быть предоставлены игровые позиции в формате, который может обрабатываться алгоритмом, позволяя модели принимать решения на основе изученных шаблонов и оценок состояния.

Всего на доске расположено $N = 24$ различных пунктов и участвует $P = 2$ игрока. Для кодирования состояния каждого пункта используется вектор $\bar{c}_{ji} = (k_1, k_2, k_3, k_4)$, где индекс $i$ обозначает номер пункта, а $j$ --- номер игрока.

Компоненты вектора $\bar{c}_{ji}$ описывают количество шашек игрока $j$ на $i$-м пункте следующим образом:
\begin{itemize}
    \item Первые три координаты $k_1$, $k_2$, $k_3$ принимают значения 0 или 1 и представляют собой бинарное кодирование количества шашек (например, одна шашка --- $(1, 0, 0)$, две --- $(1, 1, 0)$, три --- $(1, 1, 1)$ и т. д.).
    \item Четвёртая координата $k_4$ задаётся по формуле $k_4 = \frac{n - 3}{2}$, где $n \in [0, 15]$ --- фактическое количество шашек на пункте. Такой способ кодирования позволяет компактно представлять большие количества шашек и тем самым существенно снижает размерность признакового пространства, ускоряя обучение и повышая обобщающую способность модели (так как на большинстве пунктов обычно находится небольшое число шашек).
\end{itemize}

Примеры кодирования:
\begin{itemize}
    \item $(1, 0, 0, 0)$ --- одна шашка;
    \item $(1, 1, 1, 0)$ --- три шашки;
    \item $(1, 1, 1, 3)$ --- девять шашек.
\end{itemize}

Таким образом, общая размерность компоненты $\bar{c}$ равна:
\begin{equation}
    |\bar{c}| = \sum\limits_{i=1}^{N} \sum\limits_{j=1}^{P} |\bar{c}_{ji}| = 192.
\end{equation}

Дополнительно кодируются:
\begin{itemize}
    \item Число шашек на баре (то есть находящихся вне игры после взятия): $\bar{b} = (b_1, b_2)$, где $b_j = \frac{n}{2}$ --- нормализованное количество шашек игрока $j$ на баре;
    \item Число выведенных (удалённых с доски) шашек: $\bar{d} = (d_1, d_2)$, где $d_j = \frac{n}{15}$ --- нормализованное количество удалённых шашек игрока $j$;
    \item Информация о текущем ходе: $\bar{t} = (t_1, t_2)$, где $t_j \in \{0, 1\}$ указывает, чей сейчас ход.
\end{itemize}

Таким образом, итоговый вектор признаков, описывающий текущее состояние игры, имеет размерность:

\begin{equation}
    |\bar{l}_1| = |\bar{c}| + |\bar{b}| + |\bar{d}| + |\bar{t}| = 192 + 2 + 2 + 2 = 198.
\end{equation}
